\documentclass[journal]{IEEEtran}

%\usepackage[retainorgcmds]{IEEEtrantools}
%\usepackage{bibentry}  
\usepackage{xcolor,soul,framed}

\usepackage[pdftex]{graphicx}
\graphicspath{{../pdf/}{../jpeg/}}
\DeclareGraphicsExtensions{.pdf,.jpeg,.png}

\usepackage[cmex10]{amsmath}
\usepackage{amssymb}
\usepackage{amsfonts}
\newcommand{\mathbfcal}[1]{\boldsymbol{\mathcal{#1}}}

\usepackage{array}
\usepackage{mdwmath}
\usepackage{mdwtab}

\usepackage[numbers,sort&compress]{natbib}

\usepackage{eqparbox}
\usepackage{url}
\usepackage{caption}

% ---------------- ALGORITHM ----------------
\usepackage[ruled,vlined,linesnumbered]{algorithm2e}

% Define \Return used in your algorithm
\SetKw{Return}{return}

% -------------------------------------------

\usepackage{multirow}
\usepackage{booktabs}
\usepackage{float}
\usepackage{titlesec}
\usepackage{subcaption}


\begin{document}
%\bstctlcite{IEEEexample:BSTcontrol}
    \title{A Physics-Guided Traffic-Bridge Interaction Framework for Multirate Reduced-Order VBI Dynamic Response Analysis under Heterogeneous Simulated Traffic Scenarios}
\author{\footnotesize
    Nidhi Agrawal\textsuperscript{1},
  %  Rehan Nath\textsuperscript{1},
  %  Aashritha Lakshmi Mallampati\textsuperscript{1},
   % Keshav Goel\textsuperscript{1},
    Shashank Gupta\textsuperscript{1},
    Shuvendu Narayan Patel\textsuperscript{2},
    Devendra Gokul Patil\textsuperscript{3}
    
    {\small
    \textsuperscript{1}Department of Computer Science and Information Systems, Birla Institute of Technology and Science, Pilani, Pilani Campus, India. \\
    \textsuperscript{2}Department of Civil Engineering, Birla Institute of Technology and Science, Pilani, Pilani Campus, India.\\
    \textsuperscript{3}Department of Mechanical Engineering, Birla Institute of Technology and Science, Pilani, Goa Campus, India. \\
    }
     % small space before corresponding author
    {\scriptsize \textbf{Corresponding Author:} Shashank Gupta} \\[0.ex]
    \scriptsize
    \texttt{p20240429@pilani.bits-pilani.ac.in, shashank.gupta@pilani.bits-pilani.ac.in, shuvendu@pilani.bits-pilani.ac.in, devendrap@goa.bits-pilani.ac.in}
}




\maketitle

\begin{abstract}
%\boldmath
The complex dynamic interaction between stochastic traffic and long-span bridges has important implications for bridge design and safety assessment. Reliable vehicle-bridge interaction (VBI) analysis requires a numerically stable and consistent coupled formulation that captures both macroscopic traffic-flow conditions and microscopic vehicle dynamics. Existing VBI frameworks typically evaluate bridge responses under prescribed or simulated traffic conditions, while the translation of forecast traffic states into physically consistent bridge responses remains insufficiently addressed. Therefore, this study develops a physics-guided traffic-bridge interaction framework to examine how traffic-load estimation affects VBI predictions. Vehicle motion is propagated using the Intelligent Driver Model (IDM) and a Minimizing Overall Braking Induced by Lane Changes (MOBIL) rule. Then, a residual spatiotemporal correction refines the resulting physics-based traffic evolution. The estimated traffic-load field is subsequently incorporated into a reduced-order VBI formulation to evaluate the corresponding bridge dynamic response. Evaluation on 600 test windows from 37 unseen simulated traffic scenarios highlights that enforcing occupancy-load support consistency reduces the maximum predicted midspan displacement from 191.12 to 40.61 mm and acceleration from 5.281 to 0.605 m/s$^2$. These findings highlight that accurate traffic estimates alone are insufficient for reliable bridge-response prediction unless spatial load consistency is preserved.
\end{abstract}


% === KEYWORDS ====================================================================
% =================================================================================
\begin{IEEEkeywords}
Vehicle-bridge interaction, traffic-load forecasting,
physics-guided learning, spatiotemporal residual correction,
reduced-order modeling, bridge dynamic response,
occupancy-load support consistency.
\end{IEEEkeywords}

\begin{figure*}[!t]
    \centering
    \includegraphics[width=\textwidth]{Framework.png}
    \caption{Overall workflow of the proposed physics-guided traffic--bridge interaction framework.}
    \label{fig:overall_framew}
\end{figure*}

\section{Introduction}

Vehicle-bridge interaction (VBI) systems have been broadly explored to investigate the vibration response of vehicles and bridge structures under moving traffic load scenarios
\cite{han2024,yin2025,huang2026,mei2025}.
As traffic flow increases and vehicle loads become more diverse, the dynamic behaviour of the coupled system becomes increasingly complex
\cite{aloisio2023,lifeng2023}.
As a result, the loads imposed by moving vehicles and the corresponding bridge response continuously influence each other
\cite{kreslin2024,shi2024,yau2024}.
This complexity becomes more pronounced when multiple vehicles with distinct spatial positions, weights, velocities, and driving behaviours are simultaneously present on the bridge structures
\cite{yin2023,wei2025}.
Consequently, realistic representation of traffic loading and accurate vehicle-bridge coupled analysis are vital for evaluating bridge dynamic response and structural safety under actual traffic conditions
\cite{cantero2024,dong2023}.

To achieve accurate bridge responses under stochastic traffic conditions, various traffic-flow models have been introduced to acquire the randomness of realistic traffic within vehicle-bridge coupled systems
\cite{wang2024,huang2026}.
The cellular automaton (CA) model has been widely utilized for traffic simulation, since prescribed driving rules, including acceleration, deceleration, and lane-changing behaviour, can be represented in discrete space and time
\cite{rickert1996,ruan2017}.
However, the discrete-cell formulation imposes restrictions on vehicle spacing and velocity, which may limit the accurate reproduction of realistic traffic conditions
\cite{ruan2017}.
In comparison with CA-based approaches, the intelligent driver model (IDM) provides a deterministic time-space continuous representation which can describe precise vehicle positions and microscopic driving behaviour on the bridge
\cite{treiber2000,geng2023,liu2023}.
Existing IDM-based techniques have also incorporated lane-changing behaviour to simulate multi-lane traffic flow under distinct congestion scenarios
\cite{caprani2016,obrien2015,xiong2020}.

Although IDM-based traffic simulation can reproduce realistic vehicle motion, the initial spatial position, vehicle weight, and speed required for complete bridge-loading representation are not determined by the IDM alone
\cite{wangxu2019}.
In previous vehicle-bridge coupled studies, vehicle and bridge responses have often been solved as separate subsystems, with vehicle contact forces calculated in one environment and subsequently applied to the bridge model in another
\cite{yangfonder1996,green1994,pan2002}.
The bridge displacement response is then returned to the vehicle subsystem, and the interaction forces are recalculated iteratively until convergence is achieved.
Such an approach requires repeated information exchange between the vehicle and bridge models and may involve extensive coding and multiple software platforms. Therefore, recent traffic-bridge interaction frameworks have sought to integrate traffic simulation information, vehicle models, and bridge models within a unified coupled system for dynamic response analysis. Hamilton's principle has been increasingly adopted for the dynamic analysis of VBI systems, since it provides an energy-based variational framework that incorporates the kinetic and potential energies of the coupled system
\cite{mo2024}.
Through the stationary condition of the action functional, the equations of motion governing the vehicle and bridge can be systematically derived for vibration-response analysis.
Existing studies have employed Hamilton's principle to formulate coupled dynamic equations for beam-vehicle systems and other bridge-related dynamic problems involving nonlinear structural behaviour.
In recent VBI formulations, the bridge has been represented using Euler-Bernoulli beam theory, while the vehicle subsystem is described by a spring-damper model and coupled with the bridge through the relative vehicle-bridge motion.
The resulting coupled dynamic equations can subsequently be solved using the Newmark-$\beta$ method to obtain the time-dependent response of the VBI system
\cite{newmark1959,mo2024}.

The information obtained from the traffic-flow model can be incorporated into the vehicle model to determine the traffic loading applied to the bridge. In particular, vehicle information such as spatial location, velocity, weight, motion, and driver behaviour provides the basis for establishing the traffic conditions required by the traffic-bridge interaction system. The resulting traffic loads are subsequently introduced into the coupled vehicle-bridge system for dynamic response analysis. Compared with simplified moving-force or influence-line approaches, consideration of the vehicle-bridge coupled effect can provide a more realistic representation of bridge response under stochastic traffic flow
\cite{yang2005,cantero2010}.
In the coupled VBI formulation, the external excitation includes the traffic loads, while the interaction between the vehicle and bridge determines the resulting displacement and vibration response of the system.

Although the traffic state generated by the IDM can reproduce realistic vehicle motion, the initial spatial location, vehicle weight, and speed of all vehicles on the bridge are not provided by the IDM alone.
Moreover, conventional vehicle-bridge coupled analyses may require the vehicle and bridge to be solved iteratively as two independent subsystems, resulting in extensive coding and the use of multiple software applications.
On the structural side, existing VBI models are also often established using simplified assumptions, which may limit their ability to represent the complex dynamic characteristics of the coupled system. These limitations indicate the need for a framework in which realistic traffic information can be incorporated into the vehicle-bridge coupled system and subsequently utilized for bridge dynamic response analysis.


To go beyond the existing work, this study develops a physics-guided traffic-bridge interaction framework to evaluate the dynamic response of a bridge under evolving heterogeneous traffic loads.
First, the traffic information is represented in terms of vehicle spatial location, velocity, weight, lane distribution, and driving behaviour, and the subsequent traffic state is obtained using microscopic vehicle-motion models.
A residual spatiotemporal correction is then introduced to compensate for deviations in the physics-based traffic evolution and obtain the resulting traffic-load field.
Subsequently, the traffic-load information is incorporated into the vehicle model and coupled with the bridge model to establish the complete VBI system.
Finally, the bridge and vehicle subsystems are represented using the Euler-Bernoulli beam and spring-damper models, respectively, and the coupled dynamic response is acquired through Hamilton's principle and Newmark-$\beta$ time integration.



\section{Methodology}
\subsection{Overview of the Proposed Framework}
\label{sec:framework_overview}

The proposed framework addresses VBI under evolving heterogeneous traffic, where the structural response depends on the magnitude, spatial distribution, and temporal evolution of vehicle loading. Instead of treating traffic evolution and bridge dynamics independently, the framework interconnects the estimated traffic state directly to the subsequent VBI analysis through a spatially distributed traffic-load field.

Let the observed traffic history be denoted by $\mathbf{X}_{\mathrm{obs}}$, and let $\widehat{\mathbf{Y}}$ represent the estimated future traffic-load state. Hence, the overall response-estimation procedure can be expressed compactly as

\begin{equation}
\widehat{\mathbf{R}}
=
\mathcal{S}_{\mathrm{VBI}}
\left[
\mathcal{T}_{\mathrm{TB}}
\left(
\widehat{\mathbf{Y}}
\right)
\right],
\qquad
\widehat{\mathbf{Y}}
=
\mathcal{F}_{\mathrm{tr}}
\left(
\mathbf{X}_{\mathrm{obs}}
\right),
\label{eq:overall_framework}
\end{equation}

where $\mathcal{F}_{\mathrm{tr}}$ denotes the traffic-state estimation process, $\mathcal{T}_{\mathrm{TB}}$ represents the traffic-to-bridge load transformation, and $\mathcal{S}_{\mathrm{VBI}}$ is the coupled VBI solution operator.

Unlike an aggregate-load representation, the estimated traffic field retains the lane-wise and longitudinal distribution of vehicle loading. The resultant bridge response is evaluated in terms of vertical displacement and acceleration at selected locations along the structural span. The complete workflow is illustrated in Fig.~\ref{fig:overall_framew}.

\subsection{Bridge-Centric Traffic State and Load-Field Representation}
\label{sec:traffic_state_representation}

\subsubsection{Vehicle-Level Traffic State}
\label{subsec:vehicle_state}

Initially, the traffic condition is represented at the individual-vehicle level so that heterogeneous vehicle motion and loading are retained before construction of the spatial bridge-load field. Let the state of vehicle $i$ at time $t$ be expressed as

\begin{equation}
\mathbf{s}_{i}(t)
=
\left[
x_i(t),\,
\ell_i(t),\,
v_i(t),\,
w_i,\,
c_i
\right]^{T},
\label{eq:vehicle_state}
\end{equation}

where $x_i(t)$ denotes the longitudinal position, $\ell_i(t)$ indicates the lane index, $v_i(t)$ denotes the travelling speed, $w_i$ is the vehicle weight, and $c_i$ represents the vehicle-type category. Hence, the set of vehicles simultaneously present on the bridge can be expressed as

\begin{equation}
\mathcal{V}(t)
=
\left\{
\mathbf{s}_{i}(t)
\right\}_{i=1}^{N_v(t)},
\label{eq:vehicle_set}
\end{equation}

where $N_v(t)$ denotes the number of vehicles present within the traffic domain at time $t$.

This representation preserves the vehicle-level information required for constructing a spatially distributed traffic-load field for the subsequent VBI analysis.


\subsubsection{Lane-Space Discretization}
\label{subsec:lane_space}

The traffic domain is discretized in the transverse and longitudinal directions. Suppose $L_{\mathrm{tr}}$ denotes the traffic-domain length, $N_{\ell}$ the number of lanes, and $N_x$ the number of longitudinal cells. The longitudinal cell size is

\begin{equation}
\Delta x
=
\frac{L_{\mathrm{tr}}}{N_x}.
\label{eq:cell_size}
\end{equation}

For a vehicle at position $x_i(t)$, the corresponding longitudinal cell is

\begin{equation}
s_i(t)
=
\min\left\{
N_x,\,
1+\left\lfloor
\frac{x_i(t)}{\Delta x}
\right\rfloor
\right\},
\label{eq:cell_index}
\end{equation}

where the clipping at $N_x$ keeps the downstream boundary within the final one-based cell. Its transverse location is identified by $\ell_i(t)$. Hence, each vehicle is assigned to a lane-cell location

\begin{equation}
\mathcal{C}_i(t)
=
\left(
\ell_i(t),\,s_i(t)
\right).
\label{eq:lane_cell_location}
\end{equation}

In the implemented traffic representation,

\begin{equation}
N_{\ell}=4,
\qquad
N_x=50,
\label{eq:grid_configuration}
\end{equation}

so that the spatial arrangement of vehicle loading is retained before its transfer to the VBI model. In the subsequent compact notation, $L=N_{\ell}$ denotes the number of lanes and $S=N_x$ denotes the number of longitudinal cells.


\begin{figure*}[!t]
\centering

\begin{minipage}[b]{0.46\textwidth}
    \centering
    \includegraphics[width=\linewidth]{2.1.png}
\end{minipage}
\hfill
\begin{minipage}[b]{0.46\textwidth}
    \centering
    \includegraphics[width=\linewidth]{2.2.png}
\end{minipage}

\vspace{3mm}

\begin{minipage}[b]{0.46\textwidth}
    \centering
    \includegraphics[width=\linewidth]{2.3.png}
\end{minipage}
\hfill
\begin{minipage}[b]{0.46\textwidth}
    \centering
    \includegraphics[width=\linewidth]{2.4.png}
\end{minipage}

\caption{Multi-channel traffic-load matrix representation for one traffic scenario.}
\label{fig:multi_channel_traffic}
\end{figure*}


\subsubsection{Vehicle-State-to-Load-Field Mapping}
\label{subsec:load_field_mapping}

Let $\mathcal{V}_{\ell,s}(t)$ represent the set of vehicles occupying lane $\ell$ and longitudinal cell $s$ at time $t$. The local occupancy is

\begin{equation}
O_{\ell,s}(t)
=
\left|
\mathcal{V}_{\ell,s}(t)
\right|.
\label{eq:cell_occupancy}
\end{equation}

For occupancy classification and load-support enforcement, the corresponding binary occupancy state is defined as

\begin{equation}
B_{\ell,s}(t)
=
\mathbb{I}\left[O_{\ell,s}(t)>0\right],
\label{eq:binary_occupancy}
\end{equation}

where $B_{\ell,s}(t)=1$ indicates that at least one vehicle is present
in the lane-cell and $B_{\ell,s}(t)=0$ otherwise.

The mean vehicle speed in an occupied cell is
\begin{equation}
V_{\ell,s}(t)
=
\frac{1}{O_{\ell,s}(t)}
\sum_{i\in\mathcal{V}_{\ell,s}(t)}
v_i(t),
\qquad
O_{\ell,s}(t)>0,
\label{eq:mean_vehicle_speed}
\end{equation}
with
\begin{equation}
V_{\ell,s}(t)=0,
\qquad
O_{\ell,s}(t)=0,
\label{eq:empty_vehicle_speed}
\end{equation}
for empty cells.

Hence, the corresponding traffic load is defined as

\begin{equation}
P_{\ell,s}(t)
=
\sum_{i\in\mathcal{V}_{\ell,s}(t)}
w_i,
\label{eq:cell_load}
\end{equation}

so that both the magnitude and spatial location of vehicle loading are retained.

Similarly, the mean vehicle-type code is expressed as

\begin{equation}
C_{\ell,s}(t)
=
\frac{1}{O_{\ell,s}(t)}
\sum_{i\in\mathcal{V}_{\ell,s}(t)}
c_i,
\qquad
O_{\ell,s}(t)>0,
\label{eq:mean_vehicle_type}
\end{equation}

with

\begin{equation}
C_{\ell,s}(t)=0,
\qquad
O_{\ell,s}(t)=0,
\label{eq:empty_vehicle_type}
\end{equation}

for empty cells.

Therefore, the lane-cell traffic state is

\begin{equation}
\mathbf{Y}_{\ell,s}(t)
=
\left[
O_{\ell,s}(t),\,
V_{\ell,s}(t),\,
P_{\ell,s}(t),\,
C_{\ell,s}(t)
\right]^{T},
\label{eq:lane_cell_state}
\end{equation}

and the complete bridge-traffic field is gathered over all lanes and longitudinal cells as

\begin{equation}
\mathbf{Y}(t)
=
\left\{
\mathbf{Y}_{\ell,s}(t)
\right\}_{
\substack{
\ell=1,\ldots,N_{\ell}\\
s=1,\ldots,N_x
}},
\label{eq:traffic_field}
\end{equation}

with

\begin{equation}
\mathbf{Y}(t)
\in
\mathbb{R}^{N_{\ell}\times N_x\times 4}.
\label{eq:traffic_field_dimension}
\end{equation}


\subsubsection{Spatiotemporal Load-Field Assembly}
\label{subsec:spatiotemporal_field}

The lane-space representation is extended over time to form the spatiotemporal traffic tensor

\begin{equation}
\mathbfcal{Y}
=
\left\{
\mathbf{Y}(t_1),
\mathbf{Y}(t_2),
\ldots,
\mathbf{Y}(t_T)
\right\}.
\label{eq:spatiotemporal_tensor}
\end{equation}

Therefore, each traffic frame retains the lane-wise and longitudinal distributions of occupancy, mean speed, vehicle weight, and vehicle-type information.

For a reference instant $t$, the observed sequence is expressed as

\begin{equation}
\mathbf{X}_{\mathrm{obs}}
=
\left[
\mathbf{Y}_{t-T_{\mathrm{obs}}+1},
\ldots,
\mathbf{Y}_{t}
\right],
\label{eq:observed_sequence}
\end{equation}

and the subsequent traffic-load sequence by

\begin{equation}
\mathbf{Y}_{\mathrm{future}}
=
\left[
\mathbf{Y}_{t+1},
\ldots,
\mathbf{Y}_{t+T_{\mathrm{pred}}}
\right].
\label{eq:future_sequence}
\end{equation}

The implemented configuration utilizes four lanes, 50 longitudinal cells, and four traffic channels, with

\begin{equation}
T_{\mathrm{obs}}=20,
\qquad
T_{\mathrm{pred}}=30.
\label{eq:sequence_lengths}
\end{equation}

Since the traffic channels have distinct physical magnitudes, channel-wise normalization is applied as

\begin{equation}
\widetilde{Y}^{(c)}
=
\frac{Y^{(c)}}{\alpha_c},
\label{eq:channel_normalization}
\end{equation}

where $\alpha_c$ represents the scaling factor for channel $c$. The vehicle-weight channel is returned to its physical scale before being supplied to the traffic--bridge coupling procedure.

A representative multi-channel traffic-load field is shown in Fig.~\ref{fig:multi_channel_traffic}, while Fig. \ref{fig:spatiotemporal_traffic} illustrates the corresponding spatiotemporal evolution of occupancy, mean speed, and traffic loading.

\subsection{Physics-Guided Traffic Evolution Estimation}
\label{sec:physics_guided_evolution}

Primarily, the future bridge-traffic state is propagated by utilizing physically interpretable vehicle-motion models and is subsequently refined by a learned residual correction. The vehicle configuration at the final observed instant provides the initial condition for the future evolution. Longitudinal motion is described through the Intelligent Driver Model (IDM), while lane-changing behaviour is represented by utilizing a simplified MOBIL-like rule.

Let the complete vehicle configuration at the final observed instant be denoted by

\begin{equation}
\mathcal{V}_{t}
=
\left\{
z_i(t)
\right\}_{i=1}^{N_t}.
\end{equation}

Its subsequent physics-based evolution is represented by

\begin{equation}
\mathcal{V}_{t+\tau}^{\mathrm{phy}}
=
\Phi_{\mathrm{phy}}
\left(
\mathcal{V}_{t+\tau-1}^{\mathrm{phy}}
\right),
\qquad
\tau=1,\ldots,T_{\mathrm{pred}},
\end{equation}

here $\Phi_{\mathrm{phy}}$ represents the microscopic traffic-update operator. The resulting vehicle states are mapped back to the lane-space representation defined in Section~\ref{sec:traffic_state_representation}.

Fig. \ref{fig:spatiotemporal_traffic} illustrates a representative spatiotemporal traffic sequence in terms of occupancy, mean speed, and traffic-load distribution.

\begin{figure*}[!t]
\centering

\begin{minipage}[b]{0.32\textwidth}
    \centering
    \includegraphics[width=\linewidth]{3.1.png}\\
    (a)
\end{minipage}
\hfill
\begin{minipage}[b]{0.32\textwidth}
    \centering
    \includegraphics[width=\linewidth]{3.2.png}\\
    (b)
\end{minipage}
\hfill
\begin{minipage}[b]{0.32\textwidth}
    \centering
    \includegraphics[width=\linewidth]{3.3.png}\\
    (c)
\end{minipage}

\caption{Spatiotemporal bridge-traffic state representation for a representative traffic sequence: (a) vehicle occupancy, (b) mean vehicle speed, and (c) traffic-load distribution over the bridge span during the prediction interval.}
\label{fig:spatiotemporal_traffic}
\end{figure*}


\subsubsection{Longitudinal Vehicle Motion Using IDM}

For a following vehicle $i$ and its leader $j$, the longitudinal spacing can be written as:

\begin{equation}
d_i(t)
=
x_j(t)-x_i(t)-L_j,
\end{equation}

and the relative velocity is

\begin{equation}
\Delta v_i(t)
=
v_i(t)-v_j(t).
\end{equation}

The desired dynamic spacing is given by

\begin{equation}
d_i^{*}(t)
=
d_0
+
v_i(t)T_h
+
\frac{
v_i(t)\Delta v_i(t)
}{
2\sqrt{ab}
},
\end{equation}

where $d_0$ is the minimum spacing, $T_h$ the desired time headway, $a$ the maximum acceleration, and $b$ the comfortable deceleration.

The IDM acceleration is expressed as:

\begin{equation}
a_i(t)
=
a
\left[
1
-
\left(
\frac{v_i(t)}{v_0}
\right)^{\delta}
-
\left(
\frac{d_i^{*}(t)}{d_i(t)}
\right)^2
\right].
\end{equation}

The vehicle speed and position are then updated using the traffic time step $\Delta t$. This produces the longitudinal evolution required for reconstructing the future lane-space traffic field.


\subsubsection{Lane-Changing Behaviour Using a MOBIL-Like Rule}

Lane-changing behaviour is represented using a simplified MOBIL-like rule based on incentive and safety conditions. A candidate lane change is accepted when the expected acceleration benefit exceeds the prescribed threshold,

\begin{equation}
\Delta a_i
+
p
\left(
\Delta a_{\mathrm{new}}
+
\Delta a_{\mathrm{old}}
\right)
>
\Delta a_{\mathrm{th}},
\end{equation}

where $p$ is the politeness factor and $\Delta a_{\mathrm{th}}$ is the minimum incentive threshold.

The lane change must also satisfy the safety condition

\begin{equation}
a_{\mathrm{new}}^{\mathrm{after}}
\geq
-b_{\mathrm{safe}},
\end{equation}

where $b_{\mathrm{safe}}$ is the maximum admissible deceleration of the new following vehicle.

When both conditions are satisfied, the vehicle lane index is updated while its weight and vehicle category remain unchanged. The longitudinal position and speed continue to evolve according to the IDM update.


\subsubsection{Physics-Based Load-Field Rollout}

After each IDM and lane-change update, the propagated vehicle configuration is converted to the lane-space representation defined in Section~\ref{sec:traffic_state_representation}:

\begin{equation}
M_{t+\tau}^{\mathrm{phy}}
=
\mathcal{R}
\left(
\mathcal{V}_{t+\tau}^{\mathrm{phy}}
\right),
\qquad
\tau=1,\ldots,T_{\mathrm{pred}},
\end{equation}

where $\mathcal{R}$ denotes the vehicle-state-to-load-field mapping operator.

Therefore, the complete physics-based future sequence is expressed as:

\begin{equation}
Y^{\mathrm{phy}}
=
\left[
M_{t+1}^{\mathrm{phy}},
M_{t+2}^{\mathrm{phy}},
\ldots,
M_{t+T_{\mathrm{pred}}}^{\mathrm{phy}}
\right].
\end{equation}

Each frame retains occupancy, speed, vehicle-weight, and vehicle-type information over the complete lane-space grid.


\subsubsection{Residual Spatiotemporal Correction}

The physics-based rollout provides an interpretable future traffic trajectory, however it may accumulate local errors in vehicle position, lane assignment, speed, and load intensity. Hence, a learned residual correction is applied to the physics-based sequence rather than predicting the complete future traffic field independently.

Let $X_t$ denote the observed traffic history and $Y^{\mathrm{phy}}$ the physics-based rollout. The residual correction is

\begin{equation}
\Delta\widehat{Y}
=
\mathcal{F}_{\theta}
\left(
X_t,
Y^{\mathrm{phy}}
\right),
\end{equation}

and the corrected traffic field is

\begin{equation}
\widehat{Y}
=
Y^{\mathrm{phy}}
+
\Delta\widehat{Y}.
\end{equation}

The correction module preserves the lane-space organization of the traffic field and estimates local deviations from the physically propagated trajectory. The resulting output has the same temporal, lane, spatial, and channel dimensions as the physics-based sequence.


\subsubsection{Occupancy Calibration and Training Objective}
\label{sec:occupancy_training}

The lane-space traffic representation is inherently sparse, since a large proportion of spatial cells remain unoccupied at a given time instant. The count-valued channel $O_{\ell,s}$ remains part of the traffic state, whereas the binary state $B_{\ell,s}$ is used for occupancy classification and load-support enforcement. Small non-zero predictions in reference-empty cells may therefore introduce artificial vehicle-load contributions when the predicted weight field is transferred to the VBI model. To reduce such false activations, the binary occupancy support is interpreted probabilistically.

Let $z_{\tau,l,s}^{\mathrm{occ}}$ denote the occupancy logit at prediction step $\tau$, lane $l$, and longitudinal cell $s$. The corresponding occupancy probability is

\begin{equation}
p_{\tau,l,s}^{\mathrm{occ}}
=
\sigma
\left(
z_{\tau,l,s}^{\mathrm{occ}}
\right),
\end{equation}

where $\sigma(\cdot)$ denotes the sigmoid function.

A binary occupancy-support mask is then defined as

\begin{equation}
\widehat{B}_{\tau,l,s}
=
\mathbb{I}\left[
p^{\mathrm{occ}}_{\tau,l,s}
\geq \eta_{\mathrm{occ}}
\right].
\label{eq:predicted_binary_occupancy}
\end{equation}

where $\eta_{\mathrm{occ}}$ is selected using the validation partition. The predicted binary occupancy mask \(\hat B_{\tau,l,s}\) is used to suppress load-related predictions in cells classified as unoccupied so that artificial vehicle loads do not propagate into the downstream VBI analysis.
The residual correction model is trained using a composite objective designed to balance traffic-state reconstruction, consistency with the physics-based rollout, occupancy classification, and suppression of false activations. Let $\Omega$ denote the set of prediction indices over time, lane, longitudinal position, and traffic channel. The reconstruction term is written as

\begin{equation}
\mathcal{L}_{\mathrm{data}}
=
\frac{1}{|\Omega|}
\sum_{(\tau,l,s,c)\in\Omega}
\omega_c
\left(
\widehat{M}_{\tau,l,s,c}
-
M_{\tau,l,s,c}
\right)^2,
\end{equation}

where $\omega_c$ controls the relative contribution of the corresponding traffic channel.

To discourage unnecessary deviations from the physics-based evolution, a rollout-consistency term is introduced as

\begin{equation}
\mathcal{L}_{\mathrm{phy}}
=
\frac{1}{|\Omega|}
\sum_{(\tau,l,s,c)\in\Omega}
\left(
\widehat{M}_{\tau,l,s,c}
-
M_{\tau,l,s,c}^{\mathrm{phy}}
\right)^2.
\end{equation}

As occupied cells form a smaller fraction of the traffic tensor, occupancy classification is additionally constrained using a weighted binary cross-entropy term, where $N$ denotes the number of time--lane--cell occupancy entries included in the corresponding loss sum.

\begin{equation}
\mathcal{L}_{\mathrm{occ}}
=
-
\frac{1}{N}
\sum_{\tau,l,s}
\left[
\lambda_{+}
B_{\tau,l,s}
\log
p_{\tau,l,s}^{\mathrm{occ}}
+
\left(
1-B_{\tau,l,s}
\right)
\log
\left(
1-p_{\tau,l,s}^{\mathrm{occ}}
\right)
\right],
\end{equation}

where $B_{\tau,l,s}$ denotes the binary reference occupancy state defined by
Eq.~\eqref{eq:binary_occupancy}, and $\lambda_{+}$ compensates for class imbalance.

An empty-cell suppression term is further utilized to penalize non-zero predictions in reference-empty regions,

\begin{equation}
\mathcal{L}_{\mathrm{empty}}
=
\frac{1}{N}
\sum_{\tau,l,s}
\left(
1-B_{\tau,l,s}
\right)
\left(
\widehat{M}_{\tau,l,s}^{(w)}
\right)^2.
\end{equation}

Since the vehicle-weight channel provides the direct interface with the subsequent VBI analysis, the displayed training objective includes an additional load-consistency term to preserve the spatial distribution of bridge loading. Let the reference and estimated longitudinal load distributions be denoted by $P_s(t)$ and $\widehat{P}_s(t)$, respectively. The corresponding penalty is

\begin{equation}
\mathcal{L}_{\mathrm{load}}
=
\frac{1}{T_{\mathrm{pred}}S}
\sum_{\tau=1}^{T_{\mathrm{pred}}}
\sum_{s=1}^{S}
\left(
\widehat{P}_{s}(t+\tau)
-
P_{s}(t+\tau)
\right)^2.
\end{equation}

Therefore, the overall training objective is expressed as:

\begin{equation}
\mathcal{J}
=
\lambda_{\mathrm{data}}\mathcal{L}_{\mathrm{data}}
+
\lambda_{\mathrm{phy}}\mathcal{L}_{\mathrm{phy}}
+
\lambda_{\mathrm{occ}}\mathcal{L}_{\mathrm{occ}}
+
\lambda_{\mathrm{empty}}\mathcal{L}_{\mathrm{empty}}
+
\lambda_{\mathrm{load}}\mathcal{L}_{\mathrm{load}},
\end{equation}

where the coefficients $\lambda_{\mathrm{data}}$, $\lambda_{\mathrm{phy}}$, $\lambda_{\mathrm{occ}}$, $\lambda_{\mathrm{empty}}$, and $\lambda_{\mathrm{load}}$ control the relative contribution of the individual terms. Their values are selected using the validation partition and remain fixed during test evaluation.

The complete physics-guided traffic-load estimation procedure is summarized in Algorithm~\ref{alg:traffic_load_estimation}. The pseudocode propagates the vehicle set available at the end of the observation interval and does not specify a separate future-arrival mechanism.

\begin{algorithm}[t]
\caption{Physics-Guided Spatiotemporal Load-Field Estimation}
\label{alg:traffic_load_estimation}
\footnotesize

\KwData{
Observed traffic-load sequence $\mathbf{X}_{t}$,
last observed vehicle state $\mathcal{V}_{t}$,
prediction interval $T_{\mathrm{pred}}$,
trained residual mapping $\mathcal{F}_{\theta}$,
occupancy threshold $\eta_{\mathrm{occ}}$
}

\KwResult{
Estimated traffic-load field
$\widehat{\mathbf{Y}}=
[\widehat{\mathbf{M}}_{t+1},\ldots,
\widehat{\mathbf{M}}_{t+T_{\mathrm{pred}}}]$
}

$\mathcal{V}^{\mathrm{phy}}_{t} \leftarrow \mathcal{V}_{t}$\;

\For{$\tau \leftarrow 1$ \KwTo $T_{\mathrm{pred}}$}{

    \ForEach{vehicle $i \in \mathcal{V}^{\mathrm{phy}}_{t+\tau-1}$}{

        Identify the preceding vehicle $j$ in lane $l_i$\;

        $d_i \leftarrow
        x_j-x_i-L_j$\;

        $\Delta v_i \leftarrow v_i-v_j$\;

        $d_i^{*} \leftarrow
        d_0+v_iT_h+
        \dfrac{v_i\Delta v_i}{2\sqrt{ab}}$\;

        $a_i \leftarrow
        a\!\left[
        1-\left(\dfrac{v_i}{v_0}\right)^{\delta}
        -\left(\dfrac{d_i^{*}}{d_i}\right)^2
        \right]$\;

        $v_i^{+}
        \leftarrow
        v_i+a_i\Delta t$\;

        $x_i^{+}
        \leftarrow
        x_i+v_i\Delta t+
        \dfrac{1}{2}a_i(\Delta t)^2$\;

        Evaluate the MOBIL incentive condition
        \[
        \Delta a_i+
        p(\Delta a_{\mathrm{new}}
        +\Delta a_{\mathrm{old}})
        >
        \Delta a_{\mathrm{th}}
        \]\;

        Evaluate the MOBIL safety condition
        \[
        a_{\mathrm{new}}^{\mathrm{after}}
        \geq -b_{\mathrm{safe}}
        \]\;

        \eIf{both MOBIL conditions are satisfied}{
            $l_i^{+}\leftarrow l_i^{\mathrm{new}}$\;
        }{
            $l_i^{+}\leftarrow l_i$\;
        }

        $\mathbf{z}_{i}^{\mathrm{phy}}
        \leftarrow
        [x_i^{+},\,l_i^{+},\,v_i^{+},\,w_i,\,c_i]^{T}$\;
    }

    $\mathcal{V}^{\mathrm{phy}}_{t+\tau}
    \leftarrow
    \{
    \mathbf{z}_{i}^{\mathrm{phy}}
    \}_{i=1}^{N_{t+\tau}}$\;

    $\mathbf{M}^{\mathrm{phy}}_{t+\tau}
    \leftarrow
    \mathcal{R}
    \left(
    \mathcal{V}^{\mathrm{phy}}_{t+\tau}
    \right)$\;
}

$\mathbf{Y}^{\mathrm{phy}}
\leftarrow
[
\mathbf{M}^{\mathrm{phy}}_{t+1},
\ldots,
\mathbf{M}^{\mathrm{phy}}_{t+T_{\mathrm{pred}}}
]$\;

$\Delta\widehat{\mathbf{Y}}
\leftarrow
\mathcal{F}_{\theta}
\left(
\mathbf{X}_{t},
\mathbf{Y}^{\mathrm{phy}}
\right)$\;

$\widehat{\mathbf{Y}}
\leftarrow
\mathbf{Y}^{\mathrm{phy}}
+
\Delta\widehat{\mathbf{Y}}$\;

\For{$\tau \leftarrow 1$ \KwTo $T_{\mathrm{pred}}$}{
    \For{$l \leftarrow 1$ \KwTo $L$}{
        \For{$s \leftarrow 1$ \KwTo $S$}{

            $p^{\mathrm{occ}}_{\tau,l,s}
            \leftarrow
            \sigma
            \left(
            z^{\mathrm{occ}}_{\tau,l,s}
            \right)$\;

            $\widehat{B}_{\tau,l,s}
            \leftarrow
            \mathbb{I}
            \left[
            p^{\mathrm{occ}}_{\tau,l,s}
            \geq
            \eta_{\mathrm{occ}}
            \right]$\;

            \If{$\widehat{B}_{\tau,l,s}=0$}{
                Suppress load-related values in cell $(l,s)$\;
            }
        }
    }
}

$\widehat{\mathbf{Y}}
\leftarrow
\operatorname{InverseScale}
\left(
\widehat{\mathbf{Y}}
\right)$\;

\Return{$\widehat{\mathbf{Y}}$}\;

\end{algorithm}

\begin{figure}[!t]
    \centering
    \includegraphics[width=\linewidth]{4.png}
    \caption{Lane-wise spatial traffic-load distribution.}
    \label{fig:load distrib}
\end{figure}

\subsection{Traffic-Bridge Coupling and Load Transfer}
\label{sec:traffic_bridge_transfer}

The estimated traffic-load field is defined over discrete lane--space cells, whereas the VBI formulation operates in structural coordinates. Therefore, a traffic-to-bridge coupling procedure is introduced to transfer the predicted vehicle loading to the structural model while retaining its spatial distribution. The vehicle-weight channel is first returned to its physical scale and preserved separately for each lane and longitudinal cell rather than being reduced to a single total bridge load.

Let $P_{\ell,s}(t)$ denote the vehicle load associated with lane $\ell$ and longitudinal cell $s$ at time $t$. Since the traffic and structural domains may have different physical lengths, the longitudinal cell location is first expressed in normalized form,

\begin{equation}
x_s^{\mathrm{tr}}
=
\left(s-\frac{1}{2}\right)
\frac{L_{\mathrm{tr}}}{S},
\qquad
s=1,\ldots,S,
\label{eq:traffic_cell_center}
\end{equation}

and the normalized longitudinal coordinate is

\begin{equation}
\xi_s
=
\frac{x_s^{\mathrm{tr}}}{L_{\mathrm{tr}}},
\label{eq:normalized_cell_coordinate}
\end{equation}

and mapped to the structural coordinate as

\begin{equation}
x_s^{\mathrm{b}}
=
\xi_s L_{\mathrm{b}},
\end{equation}

where $L_{\mathrm{tr}}$ and $L_{\mathrm{b}}$ denote the traffic-domain and structural-span lengths, respectively. The transformation preserves relative longitudinal position rather than equal physical distance. Thus, the 500~m traffic coordinate is normalized before it is projected onto the 318~m structural span.


\subsubsection{Full-Field Spatial Load Projection}

For lane $\ell$, the instantaneous total traffic load is acquired from the contributions of all longitudinal cells,

\begin{equation}
P_{\ell}(t)
=
\sum_{s=1}^{S}
P_{\ell,s}(t).
\end{equation}

Rather than applying this load at a single equivalent position, the contribution of each occupied lane-cell is retained in the structural projection. Let $\phi_r(x)$ denote the $r$-th bridge mode shape. Then, the modal contribution of the lane-space traffic field can be written in the general form as:

\begin{equation}
Q_r^{\mathrm{tr}}(t)
=
\sum_{\ell=1}^{L}
\sum_{s=1}^{S}
P_{\ell,s}(t)
\phi_r\left(x_s^{\mathrm{b}}\right).
\end{equation}

This quantity represents the modal projection of the spatial traffic-load field and is distinguished from the final VBI force generated by the coupled vehicle and bridge dynamics. Retaining the complete lane-cell distribution allows traffic states with similar total vehicle weight but distinct longitudinal load locations to produce distinct structural excitation.

A representative lane-wise spatial load distribution is shown in Fig. \ref{fig:load distrib}.


\subsubsection{Observed-History Dynamic Initialization}

The structural state at the beginning of the future traffic interval depends on the loading that has already acted on the bridge during the observation period. Hence, the VBI system is not reinitialized from rest for each estimated traffic sequence. Instead, the common observed traffic history is first transferred through the same traffic-to-bridge coupling procedure and propagated through the VBI model.

Let the observed traffic sequence be denoted by

\begin{equation}
\mathbf{X}_{\mathrm{obs}}.
\end{equation}

The dynamic state at the end of the observation interval is represented as

\begin{equation}
\mathbf{z}_{0}
=
\mathcal{I}
\left(
\mathbf{X}_{\mathrm{obs}};
\boldsymbol{\theta}_{b},
\boldsymbol{\theta}_{v}
\right),
\end{equation}

where $\mathcal{I}(\cdot)$ denotes the dynamic initialization procedure and $\boldsymbol{\theta}_{b}$ and $\boldsymbol{\theta}_{v}$ contain the adopted bridge and vehicle parameters.

The initialized state contains the bridge modal and equivalent vehicle dynamic coordinates required to continue the coupled solution. The same state $\mathbf{z}_{0}$ is then used for all estimated future traffic-load fields. Consequently, differences among the resulting bridge responses arise from differences in the estimated traffic excitation rather than from inconsistent structural initial conditions.



\subsubsection{Multirate Traffic-to-Structural Time Transfer}

The traffic-load field and the structural response are resolved at distinct temporal scales. In the implemented framework, traffic states are available at

\begin{equation}
\Delta t_{\mathrm{tr}}
=
0.5~\mathrm{s},
\end{equation}

whereas the VBI equations are integrated using

\begin{equation}
\Delta t_{\mathrm{s}}
=
0.02~\mathrm{s}.
\end{equation}

For a structural time instant $t$ located between two consecutive traffic frames $t_k$ and $t_{k+1}$, the interpolation factor is

\begin{equation}
\alpha(t)
=
\frac{t-t_k}
{t_{k+1}-t_k}.
\end{equation}

The corresponding lane-cell load is evaluated as

\begin{equation}
P_{\ell,s}(t)
=
\left[1-\alpha(t)\right]
P_{\ell,s}(t_k)
+
\alpha(t)
P_{\ell,s}(t_{k+1}).
\end{equation}

Then, the interpolated loads are utilized to update the spatial participation terms required by the VBI solver. With the adopted time increments, each 0.5 s traffic interval is resolved using 25 structural integration increments.

Linear interpolation increases the temporal resolution of the excitation supplied to the structural integrator but does not reconstruct traffic-frequency content that is absent from the original 0.5 s traffic representation. The resulting continuously updated load field is supplied to the coupled vehicle-bridge equations described in the following section. The multirate transfer procedure is illustrated in Fig. \ref{fig:multirate}.



\begin{figure}[!t]
    \centering
    \includegraphics[width=\linewidth]{5.png}
    \caption{Multirate traffic-to-VBI time transfer.}
    \label{fig:multirate}
\end{figure}




\subsection{VBI Dynamic Formulation}
\label{sec:vbi_dynamic_formulation}

The spatially and temporally resolved traffic loading from the traffic-to-bridge interface is introduced into the VBI formulation. The
bridge and vehicle subsystems interact through the relative vertical motion
between the vehicle suspension and the bridge surface. The continuous
formulation presented below provides the mechanical basis for the reduced-order
implementation developed in Section~\ref{sec:reduced_vbi}.


\subsubsection{Bridge and Vehicle Dynamic Representation}

The bridge is represented using an Euler-Bernoulli beam. Let $w(x,t)$ denote the vertical bridge displacement at longitudinal coordinate $x$. The governing equation is

\begin{equation}
\mu_b
\frac{\partial^2 w(x,t)}{\partial t^2}
+
EI
\frac{\partial^4 w(x,t)}{\partial x^4}
=
p(x,t),
\end{equation}

where $\mu_b$ is the bridge mass per unit length and $EI$ is the flexural rigidity.

For the simply supported representation,

\begin{equation}
w(0,t)=w(L_b,t)=0,
\end{equation}

and

\begin{equation}
\frac{\partial^2 w}{\partial x^2}(0,t)
=
\frac{\partial^2 w}{\partial x^2}(L_b,t)
=
0.
\end{equation}

The vehicle subsystem is represented by a vertical spring--damper
coordinate $z_i(t)$. For vehicle $i$, the relative suspension deformation is defined as

\begin{equation}
\delta_i(t)
=
z_i(t)
-
w\!\left(x_i(t),t\right)
-
r_i(t),
\end{equation}

where $x_i(t)$ is obtained from the traffic-motion model and $r_i(t)$ denotes the optional road-profile excitation.

Usually, the corresponding interaction force is expressed as

\begin{equation}
F_i(t)
=
k_i \delta_i(t)
+
c_i \dot{\delta}_i(t).
\end{equation}

The vehicle vertical equation can then be written as

\begin{equation}
m_i \ddot{z}_i(t)
+
c_i \dot{\delta}_i(t)
+
k_i \delta_i(t)
=
0,
\end{equation}

subject to the chosen vertical coordinate convention.

For the reported numerical study, the road-profile term is set to zero,
as described in Section~\ref{sec:vbi_configuration}.

%\begin{figure*}[!t]
   % \centering
    %\includegraphics[width=\textwidth]{6.png}
 %   \caption{Residual architecture for physics-guided traffic-load prediction \& future bridge-load estimation.}
  %  \label{fig:arch}
%\end{figure*}

\subsubsection{Energy Basis of the Coupled System}

The bridge kinetic and bending strain energies are

\begin{equation}
T_b
=
\frac{1}{2}
\int_{0}^{L_b}
\mu_b
\left(
\frac{\partial w(x,t)}{\partial t}
\right)^2
\,dx,
\end{equation}

and

\begin{equation}
U_b
=
\frac{1}{2}
\int_{0}^{L_b}
EI
\left(
\frac{\partial^2 w(x,t)}{\partial x^2}
\right)^2
\,dx.
\end{equation}

For the vertical vehicle coordinate,

\begin{equation}
T_{v,i}
=
\frac{1}{2}
m_i
\dot{z}_i^{\,2}(t),
\end{equation}

while the suspension energy is

\begin{equation}
U_{s,i}
=
\frac{1}{2}
k_i
\delta_i^{\,2}(t).
\end{equation}

The total Lagrangian is

\begin{equation}
\mathcal{L}
=
T-U
=
\left(
T_b+\sum_i T_{v,i}
\right)
-
\left(
U_b+\sum_i U_{s,i}
\right),
\end{equation}

where the total kinetic and potential energies are assembled from the
bridge and vehicle contributions.

Application of Hamilton's principle,

\begin{equation}
\delta
\int_{t_1}^{t_2}
\mathcal{L}\,dt
=
0,
\end{equation}

together with the non-conservative damping contribution, yields the
coupled vehicle-bridge equations of motion.


\subsubsection{Coupled Bridge Excitation}

For multiple vehicles acting on the bridge, the continuous bridge
equation may be written as

\begin{equation}
\mu_b
\frac{\partial^2 w(x,t)}{\partial t^2}
+
EI
\frac{\partial^4 w(x,t)}{\partial x^4}
=
\sum_{i\in\mathcal{A}(t)}
F_i(t)
\delta_{\mathrm{D}}
\left(
x-x_i(t)
\right),
\end{equation}

where $\mathcal{A}(t)$ denotes the set of active vehicles and $\delta_{\mathrm{D}}(\cdot)$ is the Dirac delta function.

The continuous formulation is subsequently transformed into the reduced
modal representation described in Section~\ref{sec:reduced_vbi}, where the lane-space traffic
loading is retained during construction of the coupled excitation.

\begin{algorithm}[t]
\caption{Multirate Traffic-to-Structure VBI Response Computation}
\label{alg:vbi_response}
\scriptsize

\KwData{
Observed load field $\mathbf{P}^{obs}$,
estimated load field $\widehat{\mathbf{P}}$,
mode shapes $\phi_r(x)$,
bridge/vehicle parameters
$\Theta_b,\Theta_v$,
$\Delta t_{tr}$, $\Delta t_s$
}

\KwResult{
$w(L_b/4,t)$, $w(L_b/2,t)$,
$w(3L_b/4,t)$, $\ddot w(L_b/2,t)$
}

$N_{sub}\leftarrow \Delta t_{tr}/\Delta t_s$\;

$x_s^b\leftarrow
\left(s-\frac{1}{2}\right)L_b/S,
\quad s=1,\ldots,S$\;

$\beta\leftarrow0.25,\qquad\gamma\leftarrow0.50$\;

Initialize
$\mathbf{s}_{VBI}\leftarrow\mathbf{s}_{init}$\;

\tcp{Observed-history dynamic initialization}

\For{$k\leftarrow1$ \KwTo $T_{obs}-1$}{

    \For{$j\leftarrow1$ \KwTo $N_{sub}$}{

        $\rho\leftarrow j/N_{sub}$\;

        $\mathbf{P}^{*}\leftarrow
        (1-\rho)\mathbf{P}^{obs}_{k}
        +\rho\mathbf{P}^{obs}_{k+1}$\;

        \For{$l\leftarrow1$ \KwTo $L$}{
            $P_l^{*}\leftarrow\sum_{s=1}^{S}P^{*}_{l,s}$\;

            $\alpha_{l,s}^{*}\leftarrow
            P^{*}_{l,s}/P_l^{*},
            \quad P_l^{*}>0$\;

            $h_{l,r}^{*}\leftarrow
            \sum_{s=1}^{S}
            \alpha_{l,s}^{*}\phi_r(x_s^b)$\;
        }

        $F_r^{*}\leftarrow
        \sum_{l=1}^{L}P_l^{*}h_{l,r}^{*}$\;

        $(\mathbf{M},\mathbf{C},\mathbf{K},\mathbf{f})
        \leftarrow
        \mathcal{A}
        (\mathbf{P}^{*},\mathbf{h}^{*},\mathbf{F}^{*};
        \Theta_b,\Theta_v)$\;

        $\mathbf{s}_{VBI}\leftarrow
        \operatorname{Newmark}
        (\mathbf{M},\mathbf{C},\mathbf{K},\mathbf{f},
        \mathbf{s}_{VBI},
        \Delta t_s,\beta,\gamma)$\;
    }
}

$\mathbf{s}_{VBI}^{0}\leftarrow\mathbf{s}_{VBI}$\;

$\widehat{\mathbf{P}}_{0}
\leftarrow
\mathbf{P}^{obs}_{T_{obs}}$\;

$\mathbf{s}_{VBI}\leftarrow\mathbf{s}_{VBI}^{0}$\;

\tcp{Response under the estimated traffic-load field}

\For{$\tau\leftarrow0$ \KwTo $T_{pred}-1$}{

    \For{$j\leftarrow1$ \KwTo $N_{sub}$}{

        $\rho\leftarrow j/N_{sub}$\;

        $\widehat{\mathbf{P}}^{*}\leftarrow
        (1-\rho)\widehat{\mathbf{P}}_{\tau}
        +\rho\widehat{\mathbf{P}}_{\tau+1}$\;

        \For{$l\leftarrow1$ \KwTo $L$}{
            $\widehat P_l^{*}
            \leftarrow
            \sum_{s=1}^{S}\widehat P^{*}_{l,s}$\;

            $\widehat\alpha_{l,s}^{*}
            \leftarrow
            \widehat P^{*}_{l,s}/\widehat P_l^{*}$\;

            $\widehat h_{l,r}^{*}\leftarrow            \sum_{s=1}^{S}
            \widehat\alpha_{l,s}^{*}\phi_r(x_s^b)$\;
        }

        $\widehat F_r^{*}\leftarrow
        \sum_{l=1}^{L}
        \widehat P_l^{*}\widehat h_{l,r}^{*}$\;

        $(\mathbf{M},\mathbf{C},\mathbf{K},\mathbf{f})
        \leftarrow
        \mathcal{A}
        (\widehat{\mathbf{P}}^{*},
        \widehat{\mathbf{h}}^{*},
        \widehat{\mathbf{F}}^{*};
        \Theta_b,\Theta_v)$\;

        $\mathbf{s}_{VBI}\leftarrow
        \operatorname{Newmark}
        (\mathbf{M},\mathbf{C},\mathbf{K},\mathbf{f},
        \mathbf{s}_{VBI},
        \Delta t_s,\beta,\gamma)$\;

        Store $\{q_r(t),\ddot q_r(t)\}_{r=1}^{N_m}$\;
    }
}

\ForEach{$x_p\in
\{L_b/4,L_b/2,3L_b/4\}$}{
    $w(x_p,t)\leftarrow
    \sum_{r=1}^{N_m}\phi_r(x_p)q_r(t)$\;
}

$\ddot w(L_b/2,t)\leftarrow
\sum_{r=1}^{N_m}
\phi_r(L_b/2)\ddot q_r(t)$\;

\Return{$w(L_b/4,t),\,w(L_b/2,t),\,
w(3L_b/4,t),\,\ddot w(L_b/2,t)$}\;

\end{algorithm}

\subsection{Reduced-Order VBI Computational Formulation}
\label{sec:reduced_vbi}

The continuous vehicle-bridge
equations introduced in Section~\ref{sec:vbi_dynamic_formulation} are transformed into a reduced-order system
for numerical response computation. The bridge displacement is represented
using a finite modal basis, while the vehicle subsystem is described through
lane-equivalent vertical coordinates. The original lane-wise and longitudinal
traffic-load distribution is retained during construction of the modal
excitation so that the structural response remains sensitive to the spatial
location of the moving loads.

This formulation avoids
solving the continuous bridge equation directly at each structural time step
and provides an efficient representation for repeated VBI evaluation over the
complete traffic sequence.

\begin{figure*}[!t]
    \centering
    \includegraphics[width=\textwidth]{7.png}
    \caption{Bridge modal basis.}
    \label{fig:bridge modal}
\end{figure*}

\subsubsection{Modal Representation of the Bridge Response}

The bridge displacement
is approximated using $N_m$ retained vibration
modes,

\begin{equation}
w(x,t)
\approx
\sum_{r=1}^{N_m}
\phi_r(x)q_r(t),
\end{equation}

where $\phi_r(x)$ is the $r$-th modal basis
function and $q_r(t)$ is the corresponding
generalized coordinate.

For the adopted simply
supported basis,

\begin{equation}
\phi_r(x)
=
\sin
\left(
\frac{r\pi x}{L_b}
\right).
\end{equation}

The bridge velocity and
acceleration are obtained as

\begin{equation}
\dot{w}(x,t)
=
\sum_{r=1}^{N_m}
\phi_r(x)\dot{q}_r(t),
\end{equation}

and

\begin{equation}
\ddot{w}(x,t)
=
\sum_{r=1}^{N_m}
\phi_r(x)\ddot{q}_r(t).
\end{equation}

Projection onto the
modal basis gives

\begin{equation}
M_r\ddot{q}_r(t)
+
C_r\dot{q}_r(t)
+
K_rq_r(t)
=
Q_r(t),
\qquad
r=1,\ldots,N_m,
\end{equation}

where $M_r$, $C_r$, and $K_r$ are the generalized
modal mass, damping, and stiffness, respectively.

The total generalized modal force $Q_r(t)$ is distinguished from the
direct traffic-load projection. The latter is

\begin{equation}
Q_r^{\mathrm{tr}}(t)
=
\sum_{\ell=1}^{L}
\sum_{s=1}^{S}
P_{\ell,s}(t)
\phi_r
\left(
x_s^{\mathrm{b}}
\right).
\end{equation}

The direct projection therefore depends on the instantaneous longitudinal
distribution of traffic loading rather than only on the total bridge load.
The complete coupled-system forcing is represented through $\mathbf{f}(t)$
after vehicle--bridge interaction terms are assembled. The first five retained
modal basis functions are shown in Fig. \ref{fig:bridge modal}.

\begin{figure*}[!t]
\centering

% ==================== First row ====================

\begin{minipage}[t]{0.32\textwidth}
    \centering
    \includegraphics[width=\linewidth]{8.1.png}
\end{minipage}
\hfill
\begin{minipage}[t]{0.32\textwidth}
    \centering
    \includegraphics[width=\linewidth]{8.2.png}
\end{minipage}
\hfill
\begin{minipage}[t]{0.32\textwidth}
    \centering
    \includegraphics[width=\linewidth]{8.3.png}
\end{minipage}

\vspace{4mm}

% ==================== Second row ====================

\makebox[\textwidth][c]{%
    \begin{minipage}[t]{0.32\textwidth}
        \centering
        \includegraphics[width=\linewidth]{8.4.png}
    \end{minipage}
    \hspace{0.04\textwidth}
    \begin{minipage}[t]{0.32\textwidth}
        \centering
        \includegraphics[width=\linewidth]{8.5.png}
    \end{minipage}
}

\caption{Vehicle heterogeneity in the numerical dataset:
\textbf{(a)} Gross vehicle weight distribution;
\textbf{(b)} mean observed vehicle speed distribution;
\textbf{(c)} vehicle length distribution;
\textbf{(d)} maximum-acceleration parameter distribution;
\textbf{(e)} reaction-time parameter distribution.}

\label{fig:vehicle_heterogeneity}

\end{figure*}

\subsubsection{Coupled Vehicle-Bridge Reduced System}

The retained bridge
modal coordinates are coupled with lane-equivalent vertical vehicle
coordinates. Let

\begin{equation}
\mathbf{q}(t)
=
\begin{bmatrix}
q_1(t) &
q_2(t) &
\cdots &
q_{N_m}(t)
\end{bmatrix}^{T}
\end{equation}

denote the bridge modal
vector and

\begin{equation}
\mathbf{z}(t)
=
\begin{bmatrix}
z_1(t) &
z_2(t) &
\cdots &
z_L(t)
\end{bmatrix}^{T}
\end{equation}

the lane-equivalent vehicle coordinates. Each lane-equivalent coordinate is a reduced vertical degree of freedom used for coupling with the retained bridge modes; this representation does not establish an exact instantaneous aggregation of all individual vehicles in that lane.

The complete reduced
coordinate vector is

\begin{equation}
\mathbf{u}(t)
=
\begin{bmatrix}
\mathbf{q}(t) \\
\mathbf{z}(t)
\end{bmatrix}.
\end{equation}

For lane $\ell$, the bridge
deformation participating in the coupling is generated from the lane-specific
modal participation terms,

\begin{equation}
w_{\ell}(t)
=
\mathbf{h}_{\ell}^{T}(t)
\mathbf{q}(t),
\end{equation}

where $\mathbf{h}_{\ell}(t)$ is updated from the
interpolated lane-cell load distribution.

The relative vertical
deformation is then

\begin{equation}
\delta_{\ell}(t)
=
z_{\ell}(t)
-
\mathbf{h}_{\ell}^{T}(t)
\mathbf{q}(t),
\end{equation}

and the lane-equivalent
suspension force is written as

\begin{equation}
F_{\ell}(t)
=
k_{\ell}\delta_{\ell}(t)
+
c_{\ell}\dot{\delta}_{\ell}(t).
\end{equation}

The complete reduced
VBI system is assembled as

\begin{equation}
\mathbf{M}(t)\ddot{\mathbf{u}}(t)
+
\mathbf{C}(t)\dot{\mathbf{u}}(t)
+
\mathbf{K}(t)\mathbf{u}(t)
=
\mathbf{f}(t),
\end{equation}

where $\mathbf{M}(t)$, $\mathbf{C}(t)$, and $\mathbf{K}(t)$ are the assembled
system matrices and $\mathbf{f}(t)$ contains the
time-dependent traffic excitation.

The coupling terms are updated at each structural time step using the evolving lane-space load distribution, allowing changes in the spatial traffic configuration to influence the bridge response continuously.

\subsubsection{Newmark-$\beta$ Time Integration of the Coupled System}

The reduced coupled
system is integrated using the Newmark average-acceleration scheme. At
structural time $t_n$,

\begin{equation}
\mathbf{M}_n\ddot{\mathbf{u}}_n
+
\mathbf{C}_n\dot{\mathbf{u}}_n
+
\mathbf{K}_n\mathbf{u}_n
=
\mathbf{f}_n.
\end{equation}

The numerical
parameters are

\begin{equation}
\beta = 0.25,
\qquad
\gamma = 0.50,
\end{equation}

with structural time
step

\begin{equation}
\Delta t_s = 0.02~\mathrm{s}.
\end{equation}

At each structural
increment, the interpolated lane-cell load field is utilized to update the
modal participation terms and the corresponding system matrices and excitation
vector. The resulting state is then advanced to the next time step using the
standard Newmark update relations.


\subsubsection{Bridge Response Reconstruction}

After solving the reduced coupled system, the physical bridge
response is reconstructed from the modal coordinates. The vertical displacement
is

\begin{equation}
u_b(\xi,t)
=
\sum_{r=1}^{N_m}
\phi_r(\xi)q_r(t),
\end{equation}

and the corresponding acceleration is

\begin{equation}
\ddot{u}_b(\xi,t)
=
\sum_{r=1}^{N_m}
\phi_r(\xi)\ddot{q}_r(t).
\end{equation}

In the present study, bridge displacement is evaluated at

\begin{equation}
\xi
=
\frac{L_b}{4},
\qquad
\frac{L_b}{2},
\qquad
\frac{3L_b}{4},
\end{equation}

while acceleration is evaluated at midspan,

\begin{equation}
\ddot{u}_{\mathrm{mid}}(t)
=
\ddot{u}_b
\left(
\frac{L_b}{2},t
\right).
\end{equation}

Subsequently, the reconstructed response histories are used for comparison with the simulation-reference bridge response.




\section{Experimental Setup}

\subsection{Bridge-Traffic Scenario and Load-Field Configuration}

The numerical
investigation uses a controlled simulation-based bridge-traffic dataset
comprising 250 independent scenarios. The scenarios represent sparse, normal,
and dense traffic conditions. Each scenario contains 121 traffic states sampled
at 0.5 s intervals.

The traffic domain
covers a 500 m roadway segment with four lanes and 50 longitudinal cells,
giving a spatial resolution of 10 m per cell. Each traffic state is represented
by a $4\times50\times4$ lane-space tensor
containing vehicle occupancy, mean speed, accumulated vehicle weight, and mean
vehicle-type code.

The simulated traffic includes passenger vehicles, buses/light trucks, rigid trucks, and heavy trailer-type vehicles. Their motion follows the IDM/MOBIL-based formulation described in Section~\ref{sec:physics_guided_evolution}. The Physics-Based baseline therefore operates within the same model family as the controlled simulation environment. This benchmark should be interpreted as a model-consistent numerical evaluation rather than validation under traffic-model mismatch or field traffic. The resulting scenarios include distinct vehicle concentrations, load distributions, and lane-changing activity.

For sequence-based evaluation, each sample contains 20 observed frames (10 s) followed by 30 future frames (15 s). Fig.~\ref{fig:vehicle_heterogeneity} summarizes vehicle-level heterogeneity, while Fig.~\ref{fig:traffic_regime} shows the distribution and activity characteristics of the sparse, normal, and dense traffic regimes.


\subsection{Data Partitioning and Traffic-Estimation Configuration}

To prevent information
leakage, the 250 traffic scenarios are partitioned at the scenario level into
175 training, 38 validation, and 37 test scenarios. All temporal windows
generated from the same traffic realization remain within the same partition,
preventing strongly correlated samples from appearing in both model fitting and
evaluation.

Using the 20-frame
observation window and 30-frame estimation horizon, the final dataset contains
3000 training windows, 600 validation windows, and 600 test windows.
Channel-wise scaling follows the procedure defined in Section~\ref{sec:traffic_state_representation} and is kept
unchanged across the three partitions. Before VBI analysis, the estimated
vehicle-weight channel is converted back to its physical unit of kN.

The residual correction
model is optimized using AdamW with a hidden dimension of 96, batch size of 32,
and a maximum of 40 training epochs. Validation-based monitoring is utilized
for model selection, and the canonical split and evaluation use random seed 13.
The training objective follows the formulation introduced in Section~\ref{sec:occupancy_training},
with the selected hyperparameters fixed before evaluation on the held-out test
set.

The final dataset split
and traffic-estimation configuration are summarized in Table~\ref{tab:dataset_configuration}.

\begin{table}[!t]
\centering
\caption{Dataset, split and forecast configuration}
\label{tab:dataset_configuration}
\footnotesize
\begin{tabular}{ll}
\toprule
\textbf{Item} & \textbf{Final configuration} \\
\midrule
Traffic scenarios       & 250 \\
Frames per scenario     & 121 \\
Traffic interval        & 0.5 s \\
Traffic-grid length     & 500 m \\
Lanes                   & 4 \\
Longitudinal cells      & 50 \\
Cell length             & 10 m \\
Channels                & Occupancy, speed, vehicle weight, \\
                        & and vehicle type \\
Observation window      & 20 frames / 10 s \\
Forecast horizon        & 30 frames / 15 s \\
Training scenarios      & 175 \\
Validation scenarios    & 38 \\
Test scenarios          & 37 \\
Training windows        & 3000 \\
Validation windows      & 600 \\
Test windows            & 600 \\
Canonical split seed    & 13 \\
Batch size              & 32 \\
Hidden dimension        & 96 \\
Maximum epochs          & 40 \\
\bottomrule
\end{tabular}
\end{table}




\subsection{VBI Numerical Model and Parameter Configuration}
\label{sec:vbi_configuration}

The structural-response
evaluation uses the reduced-order VBI formulation introduced in Section~\ref{sec:reduced_vbi}.
The structural span is 318~m, and the reduced model retains the first five modes. The adopted modal frequencies are 0.758, 1.533, 2.345, 3.506, and 4.022~Hz. The sine functions define the spatial modal basis, while these frequencies are treated as adopted reduced-order modal parameters. An effective bridge mass per unit length of 30,000 kg/m is
adopted as a numerical parameter for constructing the generalized modal masses.

The bridge damping
contribution and road-surface irregularity are both set to zero in the reported
numerical study. These terms are retained in the general formulation but are
not activated here, allowing differences in bridge response to be attributed primarily
to the traffic excitation and vehicle-bridge coupling.

The vehicle subsystem
is represented using lane-equivalent vertical coordinates coupled to the bridge
modes. The adopted reduced-order equivalent vehicle parameters are 63,294~kg for mass, 1,903,172~N/m for suspension stiffness, and 7,882~N$\cdot$s/m for damping. These fixed parameters define the reduced suspension representation used in the numerical model; they are not interpreted here as instantaneous aggregates of the vehicles occupying a lane.

The traffic-load field
is available at 0.5 s intervals, whereas the VBI system is integrated using a
structural time step of 0.02 s, giving 25 structural increments per traffic
interval. Lane--cell loads are linearly interpolated between consecutive traffic
states before modal projection. Time integration is performed using the Newmark
average-acceleration scheme with $\beta=0.25$ and $\gamma=0.50$.

Before introducing the
estimated future traffic-load field, the common 20-frame observed traffic
history is propagated through the same VBI system to obtain the structural
state at the end of the observation interval. Then, this state is used as the
common initial condition for all subsequent 30-frame response evaluations.


\begin{figure*}[!t]
\centering
\begin{minipage}[b]{0.47\textwidth}
    \centering
    \includegraphics[width=\linewidth]{10.1.png}
\end{minipage}
\hfill
\begin{minipage}[b]{0.47\textwidth}
    \centering
    \includegraphics[width=\linewidth]{10.2.png}
\end{minipage}

\caption{Traffic-regime and activity distribution of the synthetic bridge-traffic
dataset. Panel (a) shows the number of sparse, normal, and dense scenarios,
while panel (b) reports the mean entered vehicles and lane-change events per
scenario for each traffic regime.}
\label{fig:traffic_regime}
\end{figure*}

\subsection{Comparative Methods and Evaluation Metrics}

All comparative methods
are evaluated using the same observed traffic history, traffic-to-bridge
transformation, reduced-order VBI model, modal basis, equivalent vehicle
parameters, and Newmark-$\beta$ integration settings. Only the estimated
traffic-load field supplied during the subsequent prediction interval is varied
among the methods.

Before evaluating the
estimated future traffic fields, the common 20-frame observed history is
propagated through the VBI model to establish the dynamic state at the end of
the observation interval. Then, the same initialized state is utilized for each
method, ensuring that differences in the reconstructed bridge responses arise
from the estimated future traffic excitation rather than from different
structural initial conditions.

\begin{table*}[!t]
\centering
\caption{Traffic-state and traffic-load estimation performance on the held-out test set.}
\label{tab:traffic_performance}
\renewcommand{\arraystretch}{1.15}
\setlength{\tabcolsep}{7pt}

\begin{tabular}{|l|c|c|c|c|c|c|}
\hline
\textbf{Method}
&
\begin{tabular}[c]{@{}c@{}}\textbf{Global}\\ \textbf{RMSE}\end{tabular}
&
\begin{tabular}[c]{@{}c@{}}\textbf{Occupied}\\ \textbf{RMSE}\end{tabular}
&
\textbf{F1}
&
\begin{tabular}[c]{@{}c@{}}\textbf{Empty}\\ \textbf{FPR}\end{tabular}
&
\begin{tabular}[c]{@{}c@{}}\textbf{Load RMSE}\\ \textbf{(kN)}\end{tabular}
&
\begin{tabular}[c]{@{}c@{}}\textbf{Total-load RMSE}\\ \textbf{(kN)}\end{tabular}
\\
\hline

Persistence
& 0.1951
& 0.4433
& 0.1451
& 0.0901
& 72.70
& 706.54
\\
\hline

Physics-Based
& \textbf{0.1310}
& 0.3361
& 0.5182
& \textbf{0.0298}
& \textbf{45.23}
& 917.54
\\
\hline

ConvLSTM
& 0.2683
& 0.2752
& 0.1886
& 1.0000
& 58.39
& 852.99
\\
\hline

Transformer
& 0.2186
& 0.3246
& 0.2115
& 0.7442
& 54.06
& 671.57
\\
\hline

Direct Prediction
& 0.2208
& 0.3240
& 0.2061
& 0.7843
& 53.60
& 576.99
\\
\hline

Residual Correction
& 0.1568
& \textbf{0.2306}
& \textbf{0.5370}
& 0.0334
& 45.30
& \textbf{512.02}
\\
\hline

\end{tabular}
\end{table*}




\subsubsection{Comparative Traffic-Estimation Methods}

The six
traffic-estimation strategies are evaluated.

\textbf{Persistence:}
The final observed traffic state is repeated throughout the 15 s estimation
interval. This provides a simple baseline without explicit future
vehicle-motion modelling.

\textbf{Physics-Based:}
The final observed vehicle configuration is propagated using the IDM-based
longitudinal-motion model and the simplified MOBIL-like lane-change rule
described in Section~\ref{sec:physics_guided_evolution}. The resulting vehicle states are subsequently mapped
to the lane--space traffic-load field without any learned correction.

\textbf{ConvLSTM:}
A convolutional recurrent baseline is deployed to estimate the future traffic
tensor directly from the observed sequence while preserving the spatial
organization of the lane-space grid.

\textbf{Transformer:}
A data-driven attention-based sequence model is utilized to estimate the
subsequent lane-space traffic field from the observed history without explicit
car-following or lane-changing constraints.

\textbf{Direct Prediction:}
This model predicts the complete future traffic field directly from the
observed sequence and does not use the Physics-Based rollout as an additive
reference.

\textbf{Residual Correction:}
The residual model uses the Physics-Based rollout as its reference trajectory
and estimates only the remaining spatiotemporal deviation. Thus, the learned
correction is applied to the physically propagated traffic field rather than
replacing it completely.

For traffic-level
evaluation, the continuous Residual Correction output is retained. Before its
vehicle-weight channel is supplied to the VBI model, an occupancy-based support
constraint is applied using the validation-selected threshold
$\eta_{\mathrm{occ}}$:

\begin{equation}
\widehat{P}^{\,\mathrm{SC}}_{\tau,l,s}
=
\widehat{B}_{\tau,l,s}\widehat{P}_{\tau,l,s}.
\label{eq:support_consistent_load}
\end{equation}

Here, $\widehat{B}_{\tau,l,s}$ is the predicted binary occupancy-support mask
defined in Eq.~\eqref{eq:predicted_binary_occupancy}.
The resulting
structural-input case is referred to as Support-Consistent Residual Correction.
It is not an independently trained model. It is the Residual Correction output after enforcing consistency between predicted occupancy and vehicle-load support.


\subsubsection{Traffic-State and Traffic-Load Evaluation Metrics}

Traffic-estimation
performance is evaluated using complementary matrix-level, occupancy-based, and
load-based metrics.

\textbf{Global normalized RMSE:}

For the metric definitions below, $N_c=4$ denotes the number of traffic channels. The global RMSE is computed over all prediction steps, lanes, longitudinal cells, and traffic channels:

\begin{equation}
\mathrm{RMSE}_{\mathrm{global}}
=
\sqrt{
\frac{1}
{T_{\mathrm{pred}}LSN_c}
\sum_{\tau=1}^{T_{\mathrm{pred}}}
\sum_{l=1}^{L}
\sum_{s=1}^{S}
\sum_{c=1}^{N_c}
\left(
\widehat{M}_{\tau,l,s,c}
-
M_{\tau,l,s,c}
\right)^2
}.
\end{equation}

In this matrix-level metric, $M$ and $\widehat{M}$ denote the channel-scaled traffic tensors defined by the normalization procedure above. This metric represents the overall traffic-tensor reconstruction error.

\textbf{Occupied-cell RMSE:}

Since the traffic field is sparse, the prediction error is also evaluated only
over reference-occupied cells:

\begin{equation}
\mathrm{RMSE}_{\mathrm{occ}}
=
\sqrt{
\frac{1}{|\Omega_{\mathrm{occ}}|}
\sum_{(\tau,l,s,c)\in\Omega_{\mathrm{occ}}}
\left(
\widehat{M}_{\tau,l,s,c}
-
M_{\tau,l,s,c}
\right)^2
},
\end{equation}

where

\begin{equation}
\Omega_{\mathrm{occ}}
=
\left\{
(\tau,l,s,c)
:
B_{\tau,l,s}=1
\right\}.
\end{equation}

This metric emphasizes
prediction quality in regions containing actual vehicle information.

\textbf{Occupancy F1-score:}

Occupancy classification performance is evaluated using

\begin{equation}
F_1
=
\frac{2TP}
{2TP+FP+FN},
\end{equation}

where $TP$, $FP$, and $FN$ denote true-positive, false-positive, and false-negative occupancy classifications, respectively. These counts compare the binary reference state $B_{\tau,l,s}$ with the predicted state $\widehat{B}_{\tau,l,s}$.

\textbf{Empty-cell false-positive rate:}

False activations within reference-empty cells are measured using the
empty-cell false-positive rate. This metric is particularly relevant for VBI
analysis, since false occupancy may introduce artificial vehicle-load
contributions.

\textbf{Lane-cell traffic-load RMSE:}

The vehicle-weight channel is first transformed back to its physical unit of
kN. The lane-cell load RMSE is then evaluated directly from the reference and
estimated vehicle-load fields:

\begin{equation}
\mathrm{RMSE}_{\mathrm{load}}
=
\sqrt{
\frac{1}
{T_{\mathrm{pred}}LS}
\sum_{\tau=1}^{T_{\mathrm{pred}}}
\sum_{l=1}^{L}
\sum_{s=1}^{S}
\left(
\widehat{P}_{\tau,l,s}
-
P_{\tau,l,s}
\right)^2
}.
\end{equation}

\textbf{Total-load RMSE:}

At
each prediction step, the total bridge traffic load is evolved by summing the
vehicle-weight field over all lanes and longitudinal cells. The corresponding
RMSE is

\begin{equation}
\mathrm{RMSE}_{\mathrm{total}}
=
\sqrt{
\frac{1}{T_{\mathrm{pred}}}
\sum_{\tau=1}^{T_{\mathrm{pred}}}
\left[
\sum_{l=1}^{L}
\sum_{s=1}^{S}
\widehat{P}_{\tau,l,s}
-
\sum_{l=1}^{L}
\sum_{s=1}^{S}
P_{\tau,l,s}
\right]^2
}.
\end{equation}

Both lane-cell and
total-load errors are reported in kN. Using both local and aggregate load
metrics is important, since similar total traffic loads can still correspond to
distinct spatial load distributions over the bridge.



\subsubsection{Bridge-Response Evaluation Metrics}

Bridge-response
accuracy is evaluated by comparing the response generated from each estimated
traffic-load field with the simulation-reference response generated using the
same VBI model.

The
vertical displacement response is evaluated at the quarter-span, midspan, and
three-quarter-span locations. For a selected bridge location $\xi$, the displacement RMSE
is

\begin{equation}
\mathrm{RMSE}_{u}(\xi)
=
\sqrt{
\frac{1}{N_t}
\sum_{n=1}^{N_t}
\left[
\widehat{u}(\xi,t_n)
-
u^{\mathrm{ref}}(\xi,t_n)
\right]^2
}.
\end{equation}

Agreement
of the displacement histories is additionally quantified using the coefficient
of determination,

\begin{equation}
R_u^2(\xi)
=
1-
\frac{
\displaystyle
\sum_{n=1}^{N_t}
\left[
u^{\mathrm{ref}}(\xi,t_n)
-
\widehat{u}(\xi,t_n)
\right]^2
}{
\displaystyle
\sum_{n=1}^{N_t}
\left[
u^{\mathrm{ref}}(\xi,t_n)
-
\overline{u}^{\,\mathrm{ref}}(\xi)
\right]^2
}.
\end{equation}

The
bridge acceleration response is evaluated at midspan. Its RMSE is

\begin{equation}
\mathrm{RMSE}_{a}
=
\sqrt{
\frac{1}{N_t}
\sum_{n=1}^{N_t}
\left[
\widehat{\ddot{u}}
\left(
\frac{L_b}{2},t_n
\right)
-
\ddot{u}^{\mathrm{ref}}
\left(
\frac{L_b}{2},t_n
\right)
\right]^2
},
\end{equation}

and
the corresponding $R^2$ value is calculated in
the same manner as for displacement.

To
assess peak structural demand, the absolute peak midspan displacement is
computed for each test sequence, and the corresponding mean absolute error is
reported as

\begin{equation}
\mathrm{MAE}_{\mathrm{peak}}
=
\frac{1}{N_{\mathrm{test}}}
\sum_{j=1}^{N_{\mathrm{test}}}
\left|
\max_{n}
\left|
\widehat{u}_{j}
\left(
\frac{L_b}{2},t_n
\right)
\right|
-
\max_{n}
\left|
u_{j}^{\mathrm{ref}}
\left(
\frac{L_b}{2},t_n
\right)
\right|
\right|.
\end{equation}

Displacement
errors are reported in mm, acceleration errors in m/s$^2$, and load-related errors
in kN.




\subsection{Numerical Verification and Computational Configuration}

To
ensure a consistent numerical comparison, the same VBI solver, bridge and
vehicle parameters, modal basis, structural time step, and Newmark-$\beta$
settings are utilized for all evaluated
traffic-load inputs. The traffic interval of 0.5 s is resolved using a
structural integration step of 0.02 s, giving 25 structural increments within
each traffic interval.

The
simulation-reference bridge response is acquired by supplying the reference
traffic-load field to the same VBI formulation utilized for all estimated
traffic cases. In addition, each method starts from the same dynamic state
obtained from the observed traffic history. Therefore, differences in the
calculated bridge responses arise from the estimated future traffic-load fields
rather than from changes in numerical settings or initial conditions.



\section{Results and Discussion}

\subsection{Traffic-State and Traffic-Load Estimation Performance}
\label{sec:traffic_results}

Table~\ref{tab:traffic_performance} summarizes the
traffic-estimation performance over the 600 held-out test windows. The
Physics-Based method achieves the lowest global RMSE of 0.1310, the lowest
empty-cell false-positive rate of 0.0298, and the lowest vehicle-weight RMSE of
45.23 kN. In contrast, Residual Correction gives the lowest occupied-cell RMSE
of 0.2306 and the highest occupancy F1-score of 0.5370. These results indicate
that the physics-based evolution provides stronger overall spatial consistency,
whereas the residual correction improves the representation of traffic states
within occupied regions.

A different ranking is
observed for aggregate traffic loading. Residual Correction achieves the lowest
total-load RMSE of 512.02 kN, compared with 706.54 kN for Persistence, 917.54
kN for the Physics-Based method, 852.99 kN for ConvLSTM, 671.57 kN for the
Transformer, and 576.99 kN for Direct Prediction. This shows that the method
providing the best global matrix reconstruction does not necessarily provide
the most accurate total traffic-load estimate.

The data-driven baselines
also show differences between local and aggregate performance. ConvLSTM
produces a relatively low occupied-cell RMSE, however a high empty-cell
false-positive rate, while the Transformer and Direct Prediction provide
comparatively low total-load errors despite less consistent occupancy support.
These differences are important for VBI analysis, since the bridge response
depends not only on the total vehicle load but also on its spatial distribution
over the span.

Fig.~\ref{fig:predicted_total_load} shows the total-load evolution, Fig.~\ref{fig:horizon_load_rmse} shows the horizon-wise error, Fig.~\ref{fig:traffic_metric_comparison} compares the traffic metrics, and Fig.~\ref{fig:spatiotemporal_load_fields} shows the spatiotemporal traffic-load distribution. Together, these
results demonstrate that traffic-estimation methods can be ranked differently
depending on whether local traffic states, occupancy support, or aggregate
loading are considered. Therefore, the effect of these differences on the
resulting bridge dynamic response is examined in the following section.



\begin{figure}[!t]
\centering
\includegraphics[width=\columnwidth]{11.png}
\caption{Predicted total traffic load over the forecast horizon for a representative held-out traffic sequence.}
\label{fig:predicted_total_load}
\end{figure}

\begin{figure}[!t]
\centering
\includegraphics[width=\columnwidth]{12.png}
\caption{Horizon-wise total traffic-load RMSE for the evaluated forecasting methods.}
\label{fig:horizon_load_rmse}
\end{figure}

\begin{figure*}[!t]
\centering

\begin{minipage}[t]{0.46\textwidth}
    \centering
    \includegraphics[width=\linewidth]{13.1.png}
\end{minipage}
\hfill
\begin{minipage}[t]{0.46\textwidth}
    \centering
    \includegraphics[width=\linewidth]{13.2.png}
\end{minipage}

\vspace{3mm}

\begin{minipage}[t]{0.46\textwidth}
    \centering
    \includegraphics[width=\linewidth]{13.3.png}
\end{minipage}
\hfill
\begin{minipage}[t]{0.46\textwidth}
    \centering
    \includegraphics[width=\linewidth]{13.4.png}
\end{minipage}

\caption{
\textbf{(a)} Global RMSE
\textbf{(b)} Occupied-cell RMSE
\textbf{(c)} Occupancy F1
\textbf{(d)} Total-load RMSE.}
\label{fig:traffic_metric_comparison}
\end{figure*}

\begin{figure*}[!t]
\centering

\begin{minipage}[t]{0.46\textwidth}
    \centering
    \includegraphics[width=\linewidth]{14.1.png}
\end{minipage}
\hfill
\begin{minipage}[t]{0.46\textwidth}
    \centering
    \includegraphics[width=\linewidth]{14.2.png}
\end{minipage}

\vspace{3mm}

\begin{minipage}[t]{0.46\textwidth}
    \centering
    \includegraphics[width=\linewidth]{14.3.png}
\end{minipage}
\hfill
\begin{minipage}[t]{0.46\textwidth}
    \centering
    \includegraphics[width=\linewidth]{14.4.png}
\end{minipage}

\caption{Spatiotemporal lane-summed traffic-load fields for a common held-out traffic
sequence: \textbf{(a)} simulation ground truth,
\textbf{(b)} physics-based traffic evolution,
\textbf{(c)} raw residual correction, and
\textbf{(d)} support-consistent residual correction.}
\label{fig:spatiotemporal_load_fields}
\end{figure*}

\begin{figure*}[!t]
\centering

\begin{minipage}[t]{0.47\textwidth}
    \centering
    \includegraphics[width=\linewidth]{15.1.png}
\end{minipage}
\hfill
\begin{minipage}[t]{0.47\textwidth}
    \centering
    \includegraphics[width=\linewidth]{15.2.png}
\end{minipage}

\caption{Representative displacement and acceleration history.}
\label{fig:representative_bridge_response}
\end{figure*}

\begin{figure*}[!t]
\centering

% ---------- First row ----------
\begin{minipage}[t]{0.32\textwidth}
    \centering
    \includegraphics[width=\linewidth]{16.1.png}
\end{minipage}
\hfill
\begin{minipage}[t]{0.32\textwidth}
    \centering
    \includegraphics[width=\linewidth]{16.2.png}
\end{minipage}
\hfill
\begin{minipage}[t]{0.32\textwidth}
    \centering
    \includegraphics[width=\linewidth]{16.3.png}
\end{minipage}

\vspace{3mm}

% ---------- Second row: centered ----------
\makebox[\textwidth][c]{%
    \begin{minipage}[t]{0.32\textwidth}
        \centering
        \includegraphics[width=\linewidth]{16.4.png}
    \end{minipage}
    \hspace{0.02\textwidth}
    \begin{minipage}[t]{0.32\textwidth}
        \centering
        \includegraphics[width=\linewidth]{16.5.png}
    \end{minipage}
}

\caption{Multi-metric comparison of the simulated bridge-response errors obtained
from the evaluated traffic-load inputs:
\textbf{(a)} displacement RMSE at $L_b/4$, $L_b/2$, and $3L_b/4$;
\textbf{(b)} midspan acceleration RMSE;
\textbf{(c)} peak-displacement MAE;
\textbf{(d)} displacement $R^2$; and
\textbf{(e)} acceleration $R^2$.}
\label{fig:bridge_response_metrics}
\end{figure*}


\begin{table*}[!t]
\centering
\caption{Bridge-response performance obtained from the evaluated traffic-load inputs.}
\label{tab:bridge_response_performance}
\scriptsize
\setlength{\tabcolsep}{4pt}
\renewcommand{\arraystretch}{1.15}
\resizebox{\textwidth}{!}{%
\begin{tabular}{lccccccc}
\toprule
\textbf{Traffic input}
& \begin{tabular}[c]{@{}c@{}}$L_b/4$ \textbf{RMSE}\\ \textbf{(mm)}\end{tabular}
& \begin{tabular}[c]{@{}c@{}}$L_b/2$ \textbf{RMSE}\\ \textbf{(mm)}\end{tabular}
& \begin{tabular}[c]{@{}c@{}}$3L_b/4$ \textbf{RMSE}\\ \textbf{(mm)}\end{tabular}
& \begin{tabular}[c]{@{}c@{}}\textbf{Acc. RMSE}\\ \textbf{(m/s$^2$)}\end{tabular}
& \begin{tabular}[c]{@{}c@{}}\textbf{Peak MAE}\\ \textbf{(mm)}\end{tabular}
& \begin{tabular}[c]{@{}c@{}}\textbf{Disp.}\\ $\mathbf{R^2}$\end{tabular}
& \begin{tabular}[c]{@{}c@{}}\textbf{Acc.}\\ $\mathbf{R^2}$\end{tabular}
\\
\midrule
Persistence
& \textbf{3.190} & 4.392 & 3.416 & 0.0281 & 4.416 & 0.571 & 0.364 \\
Physics-Based
& 3.894 & \textbf{3.427} & \textbf{1.856} & \textbf{0.0269} & \textbf{2.525} & \textbf{0.765} & \textbf{0.412} \\
ConvLSTM
& 5.391 & 7.236 & 5.238 & 0.1006 & 6.423 & -0.825 & -15.517 \\
Transformer
& 4.061 & 5.160 & 3.790 & 0.0792 & 3.517 & 0.523 & -1.838 \\
Direct Prediction
& 4.011 & 5.136 & 3.675 & 0.0869 & 4.002 & 0.517 & -2.545 \\
Support-Consistent Residual
& 4.062 & 3.727 & 2.101 & 0.0321 & 2.530 & 0.730 & 0.350 \\
\bottomrule
\end{tabular}%
}
\end{table*}

\subsection{Bridge Dynamic Response under Estimated Traffic Loads}
\label{sec:bridge_response_results}

The estimated
traffic-load fields are next evaluated through the common VBI model using
identical structural parameters, numerical settings, and observed-history
initialization. Table~\ref{tab:bridge_response_performance} summarizes the six primary structural-input cases for which the complete metric set is reported, while Fig.~\ref{fig:bridge_response_metrics} provides the corresponding multi-metric comparison. The Raw Residual case is retained separately as a diagnostic for the occupancy-load support analysis. The Physics-Based method provides the strongest overall bridge-response
agreement, with a midspan displacement RMSE of 3.427 mm, a three-quarter-span
RMSE of 1.856 mm, an acceleration RMSE of 0.0269 m/s$^2$, and a peak-displacement
MAE of 2.525 mm. Its displacement and acceleration $R^2$ values are 0.765 and
0.412, respectively.

The response accuracy
varies along the bridge span. Persistence gives the lowest quarter-span
displacement RMSE of 3.190 mm, whereas the Physics-Based method performs better
at the midspan and three-quarter-span locations. The Support-Consistent
Residual method also provides good structural agreement, with displacement RMSE
values of 4.062, 3.727, and 2.101 mm at $L_b/4$, $L_b/2$, and $3L_b/4$, respectively. Its
acceleration RMSE is 0.0321 m/s$^2$ and its peak-displacement MAE is 2.530 mm.

In contrast, the purely
data-driven methods produce larger structural errors, particularly in
acceleration. ConvLSTM gives a midspan displacement RMSE of 7.236 mm and an
acceleration RMSE of 0.1006 m/s$^2$, while the Transformer and Direct Prediction
methods give acceleration RMSE values of 0.0792 and 0.0869 m/s$^2$, respectively.
These results show that favorable traffic-estimation performance does not necessarily lead to equally accurate bridge-response estimation. The across-scenario structural-error distributions are shown in Fig.~\ref{fig:bridge_error_distribution}.

Representative displacement and acceleration histories are shown in Fig.~\ref{fig:representative_bridge_response}, while Fig.~\ref{fig:bridge_response_metrics} summarizes the overall structural comparison. The difference between
traffic-load accuracy and bridge-response fidelity is examined further in
Section~\ref{sec:load_response_fidelity}.

\begin{figure*}[!t]
\centering

% ---------- First row ----------
\begin{minipage}[t]{0.32\textwidth}
    \centering
    \includegraphics[width=\linewidth]{18.1.png}
\end{minipage}
\hfill
\begin{minipage}[t]{0.32\textwidth}
    \centering
    \includegraphics[width=\linewidth]{18.2.png}
\end{minipage}
\hfill
\begin{minipage}[t]{0.32\textwidth}
    \centering
    \includegraphics[width=\linewidth]{18.3.png}
\end{minipage}

\vspace{3mm}

% ---------- Second row: centered ----------
\makebox[\textwidth][c]{%
    \begin{minipage}[t]{0.32\textwidth}
        \centering
        \includegraphics[width=\linewidth]{18.4.png}
    \end{minipage}
    \hspace{0.02\textwidth}
    \begin{minipage}[t]{0.32\textwidth}
        \centering
        \includegraphics[width=\linewidth]{18.5.png}
    \end{minipage}
}

\caption{Variation of bridge-response errors across scenarios
\textbf{(a)} $L_b/4$ displacement error distribution;
\textbf{(b)} $L_b/2$ displacement error distribution;
\textbf{(c)} $3L_b/4$ displacement error distribution;
\textbf{(d)} midspan acceleration error distribution;
\textbf{(e)} peak-midspan-displacement error distribution.}
\label{fig:bridge_error_distribution}
\end{figure*}

\begin{figure*}[!t]
\centering

% ---------- First row ----------
\begin{minipage}[t]{0.32\textwidth}
    \centering
    \includegraphics[width=\linewidth]{19.1.png}
\end{minipage}
\hfill
\begin{minipage}[t]{0.32\textwidth}
    \centering
    \includegraphics[width=\linewidth]{19.2.png}
\end{minipage}
\hfill
\begin{minipage}[t]{0.32\textwidth}
    \centering
    \includegraphics[width=\linewidth]{19.3.png}
\end{minipage}

\vspace{3mm}

% ---------- Second row: centered ----------
\makebox[\textwidth][c]{%
    \begin{minipage}[t]{0.32\textwidth}
        \centering
        \includegraphics[width=\linewidth]{19.4.png}
    \end{minipage}
    \hspace{0.02\textwidth}
    \begin{minipage}[t]{0.32\textwidth}
        \centering
        \includegraphics[width=\linewidth]{19.5.png}
    \end{minipage}
}

\caption{Relationship between traffic-load and bridge-response errors
\textbf{(a)} Total-load error versus midspan displacement error;
\textbf{(b)} total-load error versus acceleration error;
\textbf{(c)} lane-cell load error versus displacement error;
\textbf{(d)} lane-cell load error versus acceleration error;
\textbf{(e)} total-load error versus peak-displacement error.}
\label{fig:traffic_bridge_error_relationship}
\end{figure*}

\begin{figure*}[!t]
\centering

% ---------- First row ----------
\begin{minipage}[t]{0.32\textwidth}
    \centering
    \includegraphics[width=\linewidth]{20.1.png}
\end{minipage}
\hfill
\begin{minipage}[t]{0.32\textwidth}
    \centering
    \includegraphics[width=\linewidth]{20.2.png}
\end{minipage}
\hfill
\begin{minipage}[t]{0.32\textwidth}
    \centering
    \includegraphics[width=\linewidth]{20.3.png}
\end{minipage}

\vspace{3mm}

% ---------- Second row: centered ----------
\makebox[\textwidth][c]{%
    \begin{minipage}[t]{0.32\textwidth}
        \centering
        \includegraphics[width=\linewidth]{20.4.png}
    \end{minipage}
    \hspace{0.02\textwidth}
    \begin{minipage}[t]{0.32\textwidth}
        \centering
        \includegraphics[width=\linewidth]{20.5.png}
    \end{minipage}
}

\caption{Effect of occupancy-load support enforcement
\textbf{(a)} Removed predicted load versus horizon;
\textbf{(b)} raw versus support-consistent displacement error;
\textbf{(c)} raw versus support-consistent acceleration error;
\textbf{(d)} raw versus support-consistent peak-displacement error;
\textbf{(e)} distribution of absolute peak midspan acceleration.}
\label{fig:support_consistency_effect}
\end{figure*}

\subsection{Traffic-Load Estimation versus Bridge-Response Fidelity}
\label{sec:load_response_fidelity}

The results in Sections~\ref{sec:traffic_results} and~\ref{sec:bridge_response_results} show that lower traffic-load estimation error does not necessarily
lead to more accurate bridge-response prediction. Residual Correction achieves
the lowest total-load RMSE of 512.02 kN, whereas the Physics-Based method gives
a substantially larger value of 917.54 kN. However, when the corresponding
unmodified traffic fields are introduced into the VBI model, the Physics-Based
method produces lower midspan displacement and acceleration RMSE values of
3.427 mm and 0.0269 m/s$^2$, compared with 4.219 mm and 0.0552 m/s$^2$ for the raw
residual estimate.

This difference arises since
the structural response depends not only on the total vehicle load but also on
its spatial and temporal distribution over the bridge. Hence, two traffic
estimates with similar or improved aggregate loading can produce different
bridge responses when the vehicle loads are positioned distinctly along the
span.

These results demonstrate that traffic-load accuracy and bridge-response fidelity should be evaluated separately. Fig.~\ref{fig:traffic_bridge_error_relationship} shows the corresponding relationships between traffic-load errors and structural-response errors. The influence of occupancy-load support consistency on this
relationship is examined in the following section.


\subsection{Effect of Occupancy-Load Support Consistency}
\label{sec:support_consistency_results}

The raw residual estimate
may assign non-zero vehicle loads to cells predicted as unoccupied. To reduce this inconsistency, the vehicle-load field is masked using the occupancy
threshold selected from the validation data before it is supplied to the VBI model.
This produces the Support-Consistent Residual input without retraining the
forecasting model.

Applying this constraint substantially reduces the extreme structural response, as shown in Fig.~\ref{fig:support_consistency_effect}. The maximum midspan displacement decreases from 191.12 to 40.61 mm, while the maximum midspan acceleration decreases from 5.281 to 0.605 m/s$^2$. These maximum amplitudes describe extreme predicted responses and are distinct from the sequence-averaged RMSE and MAE values reported in Table~\ref{tab:bridge_response_performance}. These results indicate that
maintaining consistency between predicted occupancy and vehicle-load location
is important when estimated traffic fields are utilized for bridge dynamic
analysis.


\section{Conclusion}

This manuscript presents
a physics-guided traffic--bridge interaction framework that links evolving
traffic-load estimation with VBI analysis. The framework combines physics-based
traffic propagation, residual spatiotemporal correction, support-consistent
load transfer, and a reduced VBI model to evaluate how traffic-estimation
errors influence bridge dynamic response.

The results show that distinct
evaluation metrics lead to distinct method rankings. The Physics-Based method
gives the strongest overall bridge-response agreement, while Residual
Correction provides the lowest total-load RMSE and the best occupied-cell and
occupancy F1 performance. However, lower traffic-load error does not always
lead to lower structural-response error. Enforcing occupancy-load support
consistency substantially reduces excessive displacement and acceleration
responses produced by the raw residual load field.

These findings indicate that traffic-load estimation for VBI should be evaluated not only in terms of aggregate prediction accuracy but also in terms of the spatial consistency of the predicted loading. The present study is based on a controlled simulation dataset and a reduced-order structural model. The traffic generator and the Physics-Based rollout also use the same IDM/MOBIL model family, so the present results should be interpreted as model-consistent numerical evidence rather than field validation. Future work should therefore examine traffic-model mismatch, explicit future-arrival modeling, more detailed bridge models, structural damping and road-surface effects, and field-based traffic and bridge-response data.

\bibliographystyle{IEEEtran}
\bibliography{Ref}


\end{document}
